%% transmon-benchmark.11 -- narrow revision of v10, 2026-10-13: the operator's
%% 2026-10-13 rulings (title, validation claim, affiliations, whiteroom) and the
%% v10 Judge follow-ups. See changelog.md, provenance.md and validation.md.
%% v1-v10 and their critic siblings are unchanged.

\documentclass{anvil-paper}
\usepackage{xurl}
\usepackage{array}

%% --- URL line-breaking: the AWS-blog and documentation references carry long
%% hyphenated URLs; allow breaks at hyphens and add stretch so the
%% bibliography stays inside the measure (render-gate overfull check).
\expandafter\def\expandafter\UrlBreaks\expandafter{\UrlBreaks\do\-}
\Urlmuskip=0mu plus 1mu
%% Compact bibliography.
\renewcommand{\bibfont}{\small}
\setlength{\bibsep}{2pt plus 1pt}

%% --- Draft-phase helper macros -------------------------------------------
%% Figure placeholder: uses the rendered figure if pub-figures has produced
%% it, otherwise a framed placeholder naming the source script.
\newcommand{\anvilfig}[3][0.98]{%
  \IfFileExists{figures/#2}%
    {\includegraphics[width=#1\linewidth]{figures/#2}}%
    {\fbox{\parbox{#1\linewidth}{\centering\vspace{5em}%
      \texttt{figures/#2}\\[0.5em]%
      {\small placeholder --- render via \texttt{figures/src/#3}}%
      \vspace{5em}}}}}

\title{A differentiable finite-element solver for superconducting-qubit
design: converged frequencies and gradient re-optimization of the
DeviceLayout.jl transmon}
\author{Robb Walters\\
  {\normalsize Independent Researcher}}
% NOTE: Crutcher Dunnavant offered coauthorship at review (operator, 2026-10-13)
\date{\today}

\begin{document}
\maketitle

\begin{abstract}
The SingleTransmon example of DeviceLayout.jl publishes the qubit and
resonator frequencies of an initial and an optimized transmon design,
computed with Palace at DeviceLayout's own second-order discretization. We
cross-check those numbers with two independently implemented
finite-element solvers, GEODE-FEM and Palace. Palace at DeviceLayout's
discretization reproduces both published pairs to within $0.7\%$, and on
shared meshes the two solvers agree on the qubit mode to $10^{-4}\%$.
Under refinement both converge to qubit frequencies $7.3\%$ and $8.7\%$
above the published $4.14$ and $3.005$\,GHz, and to resonator frequencies
$2.7\%$ above, so the published values are limited by discretization;
the finest computed levels already lie $4.1\%$ and $4.7\%$ above the
published qubit values. A static route is consistent with this: a
certified two-sided bracket on the converged island capacitance with the
sapphire tensor, $[79.56, 79.76]$\,fF, read through a two-node circuit of
junction and resonator, puts the initial qubit at $4.45$\,GHz and caps it
at $4.477$\,GHz. We then repeat DeviceLayout's
two-parameter optimization with Newton steps whose Jacobian comes from
GEODE-FEM's shape derivative of the eigenfrequencies. At DeviceLayout's
discretization it meets the published stopping criterion in 3 Palace runs,
where the derivative-free optimizer used 9. Aimed at converged frequencies,
it finds that the $3.0$/$4.0$\,GHz design needs a capacitor about $19\%$
longer than the published one ($16$--$38\%$ within the extrapolation band)
and a resonator $2.4\%$ longer. Converged frequencies are extrapolations
with stated bands, not bounds; the capacitance brackets are certified for
the meshed polygonal geometry.
\end{abstract}

\section{Introduction}
\label{sec:intro}

Superconducting-qubit layouts are designed by iterating field solves: a
designer parameterizes a geometry, simulates it, compares the extracted
frequencies or Hamiltonian parameters with targets, and adjusts. The
state-of-practice automated workflow, QDesignOptimizer
\citep{eriksson2025automated}, makes the loop explicit and steers it with
analytic models because its field solver returns observables but no
derivatives of them. Open tools now cover the whole loop.
DeviceLayout.jl \citep{devicelayout} builds the layout and its mesh, and
Palace \citep{palace}, an MFEM-based solver \citep{mfem,andrej2024high},
computes the eigenmodes. DeviceLayout's SingleTransmon example
\citep{devicelayout-blog,devicelayout-example} runs the loop end to end.
It publishes an initial transmon-plus-resonator design at $4.14$ and
$5.591$\,GHz and a design optimized to targets of $3.0$ and $4.0$\,GHz, at
$3.005$ and $3.982$\,GHz, found with a derivative-free optimizer in nine
Palace runs.

Those frequencies come from one discretization, and the example itself
warns that they are sensitive to it. This paper measures how sensitive.
Our main result is that the published frequencies are
discretization-limited: two independently implemented solvers converge to
qubit frequencies $7$--$9\%$ above them, a certified bound on the island
capacitance, read through a circuit model, is consistent with that shift
and caps the initial qubit, and a gradient-assisted re-optimization at
converged resolution estimates that the optimized design needs a capacitor
about $19\%$ longer (an extrapolated estimate, $16$--$38\%$).

The second solver is GEODE-FEM (``geode-fem''), a Rust finite-element code
on the Burn tensor framework \citep{burn}. It solves the same
$H(\mathrm{curl})$ eigenproblem as Palace with an independent
implementation, and it computes design derivatives: shape and material
sensitivities of its electrostatic and driven operators, and the
Hellmann--Feynman shape derivative of an eigenfrequency. On the
committed DeviceLayout mesh, given as the same input file to both codes,
geode-fem reproduces Palace's six eigenmodes above $4.5$\,GHz to
$\leq 0.033\%$ and the qubit mode to $5\times10^{-5}\,\%$
(Section~\ref{sec:agreement}).
This parity is \emph{a credential, not the contribution} --- matching
Palace is the price of entry --- but it also gives an independent
same-mesh cross-check of the Palace/DeviceLayout.jl transmon stack,
complementing the earlier checks of Palace by SQDMetal
\citep{sommers2025open} and \citet{ye2025electromagnetic}.

Our contributions are:
\begin{itemize}
  \item \textbf{A converged cross-check of a published design example.}
    Both published geometries are meshed by DeviceLayout itself and solved
    by both codes on a resolution ladder with first- and second-order
    elements. Palace at DeviceLayout's discretization reproduces the
    published pairs to under $0.7\%$; the converged estimates lie
    $7.3\%$ and $8.7\%$ above the published qubit frequencies and $2.7\%$
    above the resonators, outside every extrapolation band
    (Section~\ref{sec:crosscheck}).
  \item \textbf{Certified capacitance brackets with the sapphire tensor.}
    A conforming upper bound and an equilibrated-flux lower bound, a
    classical pairing applied here to an island capacitance with a
    floating feedline, bracket the converged island capacitance of the
    unmodified device in $[79.56, 79.76]$\,fF. Read through the junction
    and resonator circuit they give the initial qubit at $4.45$\,GHz,
    inside the eigenmode band, and a cap of $4.477$\,GHz. The certified
    difference between the sapphire tensor and its scalar average is
    $2.5$--$3.1\%$ (Section~\ref{sec:static}).
  \item \textbf{A gradient-assisted re-optimization.} Newton steps in
    DeviceLayout's own variables, with a Jacobian from geode-fem's
    eigenfrequency shape derivative and a transmission-line model,
    meet DeviceLayout's criterion in 3 Palace runs against the published
    9, and at converged resolution move the optimized capacitor from
    $1217.678$ to about $1446\,\mu$m, an extrapolated estimate whose band
    spans $16$--$38\%$ over the published length
    (Section~\ref{sec:reopt}).
  \item \textbf{A validated sensitivity layer.} Material and geometry
    derivatives of the electrostatic and driven $H(\mathrm{curl})$
    operators, each checked against finite differences of the full
    pipeline with mutation tests (Section~\ref{sec:adjoint}).
  \item \textbf{The correctness credential.} Same-input-mesh agreement
    with Palace on six eigenmodes and the qubit mode, with the
    $17.49$\,GHz mode both codes share identified as an artifact of the
    distributed port model. It doubles as an independent same-mesh
    cross-check of the Palace/DeviceLayout.jl transmon stack, complementing
    the earlier checks of Palace by SQDMetal \citep{sommers2025open} and
    \citet{ye2025electromagnetic}.
\end{itemize}

Section~\ref{sec:related} places the work. Section~\ref{sec:setup}
describes the solvers, the device and the models;
Section~\ref{sec:agreement} the same-mesh agreement;
Sections~\ref{sec:crosscheck}--\ref{sec:reopt} the cross-check, the
certified capacitance and the re-optimization;
Section~\ref{sec:adjoint} the sensitivity layer;
Section~\ref{sec:built} how geode-fem was designed and built; and
Section~\ref{sec:discussion} what is and is not established.
Appendix~\ref{sec:appendix-supporting} reports an electrostatic design
loop with a certified acceptance test, and
Appendix~\ref{sec:appendix-perf} the performance context.

\section{Related Work}
\label{sec:related}

\paragraph{Design sensitivities in finite-element electromagnetics.}
Adjoint and direct sensitivity analysis of frequency-domain
electromagnetic solvers is an established field. \citet{nikolova2004adjoint}
review adjoint-variable methods for high-frequency structures across FEM,
MoM, TLM and FDTD. On tetrahedral vector elements, Webb and coworkers
derived design sensitivities with high-order vector elements and of the
frequency response, and used them to optimize microwave devices
\citep{webb2001design,webb2002design,nair2003optimization}, and continuum
shape sensitivity for high-frequency devices has been formulated as well
\citep{choi2011generalized}. Shape adjoints for Maxwell cavity eigenvalues
were developed for accelerator cavities, for shape optimization
\citep{akcelik2005adjoint} and for inverse shape determination of
deformed cavities \citep{akcelik2008shape}, and a recent study
optimizes Maxwell eigenvalues by shape variations with hand-derived
adjoint gradients \citep{herter2026numerical}. Commercial solvers expose
derivatives: HFSS computes analytic S-parameter derivatives with respect
to geometric and material variables by reusing the forward LU factors
\citep{ansys_hfss_derivatives,vardapetyan2008sensitivity}, and COMSOL
offers forward and adjoint sensitivities and electromagnetic shape
optimization \citep{comsol_sensitivity_method,jensen2022comsolshape}.
Automated shape differentiation for general FEM exists in UFL/Firedrake
\citep{ham2019automated} and NGSolve \citep{gangl2021fully}. Palace exposes
no derivatives. For capacitance, adjoint sensitivities are standard in
interconnect extraction \citep{elmoselhy2008efficient} and FEM
\citep{ren2012computation}, and the result that the shape derivative of a
stored-energy capacitance reduces to its explicit term is the classical
self-adjoint case \citep{park2019continuum}. What geode-fem adds is an
open implementation in which these derivatives come from the same element
kernels as the solve, applied here to a qubit design loop; the methods
are not new.

\paragraph{Superconducting-qubit design optimization.}
QDesignOptimizer \citep{eriksson2025automated} wraps
non-differentiable simulation with physics-guided updates. The SQcircuit
line \citep{rajabzadeh2023analysis,rajabzadeh2024general} and SuperGrad
\citep{wang2025supergrad} differentiate at the circuit-Hamiltonian level,
so the map from geometry to circuit parameters still runs through field
simulation. SQuADDS \citep{shanto2024squadds} replaces simulation by
lookup in a validated design database, and \citet{eun2023shape}
shape-optimize transmon pads with FEM and a derivative-free global
optimizer. Qiskit Metal's quasi-lumped methods \citep{minev2021circuit} and
energy-participation and black-box quantization
\citep{koch2007,nigg2012,minev2021} supply the parameter extraction. To our
knowledge, within qubit EM design no prior work drives a layout with
gradients from the field solve.

\paragraph{Certified two-sided capacitance bounds.}
Bracketing a capacitance between complementary energy principles is
classical, from the analytic variational bounds of
\citet{polya1951isoperimetric} to finite-element dual formulations
\citep{hammond1976calculation,penman1983dual}, and in three dimensions
\citet{specogna2011complementary} gives complementary formulations that
bound capacitance from both sides. Our lower bound uses the hypercircle
idea of \citet{prager1947approximations} with Raviart--Thomas fluxes
\citep{raviart1977mixed} equilibrated patch by patch
\citep{braess2007equilibrated,braess2009equilibrated,ern2015polynomial},
the route of guaranteed goal-oriented bounds such as
\citet{pares2021posteriori}. The method is not new. What we add is its
application to an island capacitance with a floating conductor on a
production qubit layout, used as a design-acceptance test; we found no
certified capacitance bound in the qubit-design literature.

\paragraph{Cross-checks of Palace.}
SQDMetal \citep{sommers2025open} cross-checks a Palace workflow against
COMSOL and Ansys with mesh-convergence studies, and
\citet{ye2025electromagnetic} compare a Palace workflow with cryogenic
measurement at the ${\sim}0.3\%$ level. Both use models built
independently for each tool. Our comparison instead holds the input mesh
and the junction model fixed and changes the solver implementation, and it
adds a converged reference for a published design. In V\&V terms
\citep{roache1997quantification,oberkampf2010verification} it is code
verification and code-to-code comparison, not validation against
measurement. Code-to-code benchmarking on precisely specified problems is
a long tradition in computational electromagnetics
\citep{hoffmann2009comparison}. Palace and DeviceLayout.jl have no methods
publication \citep{palace,palace-blog,devicelayout,devicelayout-blog,devicelayout-juliacon}.

\paragraph{Differentiable FEM on general tensor stacks.}
Differentiable electromagnetics is most developed in photonics, on
structured grids and integral methods
\citep{molesky2018inverse,hughes2019forward,hammond2022high,schubert2022inverse,mahlau2026fdtdx,ponomareva2025torchgdm,kim2024meent}.
Differentiable FEM on autodiff tensor stacks exists outside EM
\citep{xue2023jax,liang2023pytorch} and for general Galerkin assembly
\citep{wen2026learning,chi2026torch}. FEM kernel generation is a long
tradition \citep{kirby2006compiler,alnaes2014unified,rathgeber2017firedrake,homolya2018tsfc},
and its GPU form is libCEED \citep{brown2021libceed}, on which Palace
runs. geode-fem phrases assembly as batched tensor algebra on a
general-purpose ML stack \citep{ragankelley2013halide,chen2018tvm} so
that element-level derivatives are available by construction.

\section{Solvers, Device and Models}
\label{sec:setup}

\subsection{geode-fem}
\label{sec:architecture}

geode-fem assembles first-order N\'ed\'elec \citep{nedelec1980} and
scalar P1/P2 element operators as batched $[n_{\mathrm{elem}},k,k]$ tensor
operations on Burn, whose cubecl compiler targets CPU, CUDA and WebGPU
from one source. Its eigensolver is shift-invert Lanczos \citep{arpack} on
a real symmetric pencil, factorized with faer \citep{kazdadi2026faer};
large electrostatic solves use conjugate gradients with an
algebraic-multigrid preconditioner in native f64. Autodiff reaches
everything up to the linear solve. The sparse factorization is outside
the tape, so a discrete-adjoint layer supplies the missing derivative with
one transpose solve on the same factorization
(Section~\ref{sec:adjoint}). The eigenfrequency derivative used in
Section~\ref{sec:reopt} needs no adjoint solve, because the pencil is
symmetric: for a simple eigenpair $(\lambda, x)$ with $x^{\mathsf T}Mx =
1$,
\begin{equation}
  \frac{d\lambda}{d\theta} \;=\;
  x^{\mathsf T}\Bigl(\frac{\partial K}{\partial\theta}
  - \lambda\,\frac{\partial M}{\partial\theta}\Bigr)x ,
  \label{eq:hf-eig}
\end{equation}
the Hellmann--Feynman (Nelson) formula \citep{nelson1976simplified}. geode-fem
evaluates the element derivatives $\partial K/\partial\theta$ and
$\partial M/\partial\theta$ by forward-mode dual numbers through the same
tensor-permittivity element kernels it assembles, along a node velocity
$\partial X/\partial\theta$.

\subsection{Device, mesh and materials}
\label{sec:geometry}

The device is DeviceLayout.jl v1.15.0's SingleTransmon example
\citep{devicelayout}: a transmon with a meander readout resonator and a
feedline on sapphire, in a vacuum box (Figure~\ref{fig:geometry}). Its two
design parameters are the transmon's \texttt{capacitor\_length} $c$ and the
resonator's \texttt{total\_length} $L$; the published initial design has
$c = 620\,\mu$m and $L = 5000\,\mu$m. The sapphire permittivity is
$\mathrm{diag}(9.3, 9.3, 11.5)$ with the optic axis normal to the chip.
The configuration rotates the material axes about the chip normal by
about $36.87^\circ$; for a uniaxial tensor whose optic axis is that normal
the rotation is the identity, so the lab tensor is diagonal. Where a
scalar model is used it is the trace average $\bar\varepsilon = 10.03$.

Two mesh families appear. The committed fixture is DeviceLayout's mesh of
the initial design (22{,}684 nodes, 133{,}314 tetrahedra), pinned by
SHA-256; Section~\ref{sec:agreement} and the capacitance brackets of
Section~\ref{sec:static} use it and its refinements, which leave its
polygonal surfaces unchanged. For the cross-check and the
re-optimization, DeviceLayout itself meshes each geometry with the two
parameters exposed, at first order for the ladders and at its own second
order for the runs that reproduce its discretization. On the initial
geometry the two families give the same Palace qubit at $h$ ($3.4527$
against $3.4520$\,GHz). PEC is applied on the metal sheets and the box;
the readout ports are left open in every eigenproblem except the
DeviceLayout-configuration runs of Section~\ref{sec:crosscheck}.

\begin{figure}[tbp]
  \centering
  \anvilfig{fig1-geometry.pdf}{fig1\_geometry.py}
  \caption{Benchmark geometry, drawn from the committed SHA-256-pinned
  MSH~4.1 input mesh of the DeviceLayout.jl v1.15.0
  \texttt{SingleTransmon} device. (a) Top view of the three node-disjoint
  metal conductors of Section~\ref{sec:static}: the feedline (teal), the
  qubit island (blue), and the ground plane together with the meander
  readout resonator (gray), whose etched gaps are drawn as dark
  outlines. Squares mark the two readout ports; the dashed line is the
  exterior-boundary outline. (b) The $1\times1\,\mu$m junction
  \texttt{lumped\_element} sheet with its four-triangle triangulation.
  (c) Material cross-section, schematic and not to scale; the
  \texttt{substrate} and \texttt{vacuum} volume groups fill the enclosing
  box. Both solvers read this file (22{,}684 nodes; 133{,}314
  tetrahedra); Palace then refines it internally around the metal sheets
  (Section~\ref{sec:agreement}).}
  \label{fig:geometry}
\end{figure}

\subsection{Junction and circuit models}
\label{sec:junction}

The Josephson junction is a linear lumped element, $L_J = 14.860$\,nH in
parallel with $C_J = 5.5$\,fF, on the four-triangle sheet $\Gamma$ of
Figure~\ref{fig:geometry}b, exactly as in DeviceLayout's Palace
configuration. With port length $\ell$ and width $w$ the uniform
lumped-port construction adds
\begin{equation}
  K_{\mathrm{port}} = \frac{\ell}{w\,L_J}\, S_\Gamma
  \qquad\text{and}\qquad
  M_{\mathrm{port}} = \frac{C_J\,\ell}{w}\, S_\Gamma
  \label{eq:port}
\end{equation}
to the pencil $(K + K_{\mathrm{port}})x = \omega^2(M + M_{\mathrm{port}})x$,
where $S_\Gamma$ is the isotropic surface matrix on $\Gamma$. Palace's
\texttt{LumpedPort} adds the same terms with the same scalings, so the
default models of the two codes coincide. The isotropic sheet also stores
energy in a zero-voltage field pattern confined to the patch, which
produces a mode at $17.49$\,GHz in both codes
(Section~\ref{sec:agreement}). geode-fem's rank-1 port-voltage shunt
replaces $S_\Gamma$ by $b\,b^{\mathsf T}$, with $b_i = (1/w)\int_\Gamma
N_i\cdot\hat e\,dS$ the port-voltage functional. It removes that mode and
keeps the others to four decimals on the fixture, and it is the junction
model of every geode-fem eigensolve in Sections~\ref{sec:crosscheck}
and~\ref{sec:reopt}.

The electrostatic model splits the metal into three conductors: ground
with the resonator, the floating feedline, and the island. Its island
capacitance with the feedline floating and neutral,
$C_\Sigma = C_{ii} - C_{if}^2/C_{ff}$, excludes $C_J$. A two-node circuit
connects it to the eigenproblem: a bare qubit
$1/(2\pi\sqrt{(L_J + L_s)(C_\Sigma + C_J)})$, with $L_s$ the series
inductance of the junction's current path, coupled with coefficient $k$
to the resonator. On the fixture $L_s$ and $k$ were measured on three and
two discretizations; extrapolated, they give bands of
$0.113$--$0.125$\,nH and $0.0227$--$0.0310$, which are estimates, not
bounds.

\section{Same-Input-Mesh Agreement With Palace}
\label{sec:agreement}

Both codes read the committed fixture with the same materials and
junction model at first order. Table~\ref{tab:agreement} gives the
result. Palace returns the qubit mode only when its search target lies
below it, because its eigensolver searches above the target (read from
Palace's source and configuration schema); the six
modes above $4.5$\,GHz come from a run with \texttt{Target}~$=4.5$\,GHz
and the qubit from one with \texttt{Target}~$=2$\,GHz. All six agree to
$\leq 0.033\%$ (worst case $0.032\%$), and the qubit mode to
$5\times10^{-5}\,\%$, with port participation $0.9942$ in geode-fem and
port EPR $0.99424$ in Palace. Doubling $L_J$ moves the qubit by a ratio of
$0.70810$ in Palace and $0.7081$ in geode-fem, close to $1/\sqrt2$, and a
Palace rerun from the committed configuration reproduces its
\texttt{eig.csv} bit for bit.

\begin{table}[tbp]
  \centering
  \caption{Eigenmode agreement from the identical input mesh (22{,}684
  nodes, 133{,}314 tetrahedra), first order in both solvers, junction as
  the lumped $L_J \parallel C_J$ shunt. geode-fem solves $156{,}863$
  N\'ed\'elec unknowns ($133{,}108$ interior); Palace refines the mesh
  internally around the metal sheets and solves $183{,}891$. geode-fem
  frequencies are the committed three decimals (six for the qubit), and
  $\Delta$ is the committed relative difference. These are same-mesh
  values, not converged frequencies. Source:
  \texttt{benchmarks/transmon\_eigen/results.toml}.}
  \label{tab:agreement}
  \small
  \begin{tabular}{@{}llccc@{}}
    \toprule
    Mode & & geode-fem (GHz) & Palace (GHz) & $\Delta$ (\%) \\
    \midrule
    qubit & transmon qubit & 3.452751 & 3.452749306805 & 0.00005 \\
    resonator & readout resonator & 5.153 & 5.151335830348 & 0.032 \\
    mode 2 & box/cavity & 15.465 & 15.46052107794 & 0.029 \\
    port sheet & port-sheet artifact & 17.490 & 17.49010903536 & 0.001 \\
    mode 4 & box/cavity & 18.693 & 18.69165792915 & 0.007 \\
    mode 5 & box/cavity & 20.703 & 20.69755679425 & 0.026 \\
    mode 6 & box/cavity & 26.088 & 26.08089940472 & 0.027 \\
    \bottomrule
  \end{tabular}
\end{table}

The $17.49$\,GHz mode needs care, and its explanation below is an
inference that a synthetic test supports. In geode-fem all its stiffness energy is
in the port term ($p = 1.000$), while Palace, which builds its port EPR
from the port voltage, gives it $2.49\times10^{-8}$; its energy sits in a
surface pattern with zero net voltage. The junction patch has one free
node, whose hat-function gradient carries no port voltage but still
stores energy in the isotropic sheet term, with Rayleigh quotient
$1/(L_JC_J)$, or $f_{LC} = 17.60$\,GHz, lowered to $17.49$\,GHz by the
adjacent volume mass. A two-terminal junction cannot carry that pattern.
Its $0.001\%$ agreement therefore tests that the two codes assemble the
same port operator and nothing more, and the rank-1 shunt of
Section~\ref{sec:junction} removes the mode.

The fixture is coarse. Palace's own refinements of it move the qubit from
$3.4527$\,GHz to $3.9528$\,GHz at $h/2$ ($+14.5\%$) and to $4.1721$\,GHz
at second order on $h$ ($+20.8\%$), and the resonator by $+6.1\%$ at
$h/2$. The agreement above is code verification on one discretization;
Section~\ref{sec:crosscheck} supplies the converged numbers.

\section{The DeviceLayout Cross-Check}
\label{sec:crosscheck}

\paragraph{The published designs.}
The DeviceLayout.jl blog \citep{devicelayout-blog} reports the initial
design ($c = 620\,\mu$m, $L = 5000\,\mu$m) at $[4.14, 5.591]$\,GHz and
the final design ($c = 1217.678\,\mu$m, $L = 7243.955\,\mu$m) at
$[3.005, 3.982]$\,GHz, reached by the derivative-free PRIMA optimizer in
nine Palace runs ``within 1\% of our target frequencies''. The example
source \citep{devicelayout-example} sets the targets to $[3.0, 4.0]$\,GHz,
works in the variables $x_1 = \sqrt{620\,\mu\mathrm{m}/c}$ and $x_2 =
5000\,\mu\mathrm{m}/L$, and stops when $\mathrm{err} = \frac12\sum_m
(f_m - t_m)^2/t_m^2 < 10^{-4}$. Its frequencies are computed with Palace
on DeviceLayout's mesh at mesh order 2 and solver order 2, with a
first-order absorbing far field and two $50\,\Omega$ readout ports.

\paragraph{Reproduction at DeviceLayout's discretization.}
Palace, run on DeviceLayout's second-order mesh with the example's
configuration file, gives $[4.1683, 5.6132]$\,GHz for the initial design
and $[3.0164, 3.9956]$\,GHz for the final one, $0.68\%$ and $0.40\%$
above the published initial pair and $0.38\%$ and $0.34\%$ above the
final pair. The same cells with the PEC-box model used everywhere else
differ from the configuration-file runs by at most $0.003\%$, so the
boundary model does not matter here.

\paragraph{The ladder.}
For each geometry, Palace runs at order 1 on DeviceLayout's first-order
mesh $h$ and its uniform refinements $h/2$ and $h/4$ (up to
$12.6$ million unknowns), and at order 2 on $h$ and $h/2$; geode-fem runs
at order 1 on $h$ and $h/2$, on the same meshes. Every series rises
under refinement (Figure~\ref{fig:crosscheck}). The sheet-edge
singularity of the field, $E \sim r^{-1/2}$, limits the asymptotic
eigenvalue rate on uniform meshes to $O(h)$ at any element order; the
order-1 ladders show pre-asymptotic orders of $0.61$ (qubit) and
$0.70$--$0.73$ (resonator).

\paragraph{Extrapolation rule.}
The \emph{estimate} is the mean of two extrapolations at the asymptotic
order 1 in the frequency, $2f_2 - f_1$: one from the order-1 pair
($h/2$, $h/4$) and one from the order-2 pair ($h$, $h/2$). The
\emph{band} spans six values: those two; the same pairs extrapolated at
order 1 in the eigenvalue variable $\lambda\propto f^2$,
$\sqrt{2f_2^2 - f_1^2}$, which set the low side; and the same pairs at the
observed order, which set the high side. The band is an extrapolation
range, not a bound, and its high side is kept wide on purpose. For the
initial qubit, the static route of Section~\ref{sec:static} caps the
band at $4.477$\,GHz. The rule was fixed after the static route had been
computed: extrapolation at the observed order alone, the rule used
before, puts the initial qubit at $4.63$\,GHz, above the certified cap,
and the observed-order values are kept only as the band's high side. The
agreement of the two routes is therefore a consistency check, not a blind
prediction.

\begin{figure}[tbp]
  \centering
  \anvilfig{fig2-crosscheck.pdf}{fig2\_crosscheck.py}
  \caption{Resolution ladders of the two published DeviceLayout.jl
  SingleTransmon designs. Palace at order 1 on DeviceLayout's mesh $h$
  and its uniform refinements, Palace at order 2, and geode-fem at order 1
  on the same meshes (crosses, on top of Palace's order-1 points); the
  open diamond is Palace on DeviceLayout's own second-order mesh with its
  configuration file, the discretization of the published values (dashed).
  At the right, the converged estimate with its extrapolation band, which
  is not a bound. In (a), the static route of Section~\ref{sec:static}
  (certified capacitance bracket read through the junction--resonator
  circuit) and the circuit cap from the certified lower bound. Source:
  \texttt{benchmarks/transmon\_devicelayout\_crosscheck/results.toml}.}
  \label{fig:crosscheck}
\end{figure}

\begin{table}[tbp]
  \centering
  \caption{Published frequencies, Palace at DeviceLayout's discretization,
  and the converged estimates with their extrapolation bands, in GHz.
  Percentages are relative to the published value. The estimates and bands
  are extrapolations, not bounds. Source:
  \texttt{benchmarks/transmon\_devicelayout\_crosscheck/results.toml}.}
  \label{tab:crosscheck}
  \small
  \begin{tabular}{@{}lcccc@{}}
    \toprule
    & Published & At DeviceLayout's discretization & Converged estimate & Band \\
    \midrule
    Initial qubit & 4.14 & 4.1683 ($+0.68\%$) & 4.444 ($+7.3\%$) & $[4.415, 4.677]^{\ast}$ \\
    Initial resonator & 5.591 & 5.6132 ($+0.40\%$) & 5.740 ($+2.7\%$) & $[5.723, 5.817]$ \\
    Final qubit & 3.005 & 3.0164 ($+0.38\%$) & 3.265 ($+8.7\%$) & $[3.236, 3.471]$ \\
    Final resonator & 3.982 & 3.9956 ($+0.34\%$) & 4.088 ($+2.7\%$) & $[4.077, 4.139]$ \\
    \bottomrule
  \end{tabular}

  \smallskip
  {\footnotesize $^{\ast}$Capped at $4.477$\,GHz by the static route of
  Section~\ref{sec:static}: $[4.415, 4.477]$.}
\end{table}

\paragraph{Result.}
Table~\ref{tab:crosscheck} gives the converged estimates. The initial
qubit converges to $4.444$\,GHz and the final qubit to $3.265$\,GHz,
$7.3\%$ and $8.7\%$ above the published values; both resonators converge
$2.7\%$ above theirs. Every published value lies below its band. Measured
against the estimates, DeviceLayout's discretization sits $6.2\%$ and
$7.6\%$ low on the two qubits and $2.2\%$ and $2.3\%$ low on the
resonators, and the first-order mesh $h$ alone sits $22.3\%$ and $26.7\%$
low on the qubits.

\paragraph{Same-mesh agreement on the ladder.}
geode-fem and Palace agree on the qubit to $10^{-4}\,\%$ at every shared
level of both geometries. On the resonator they differ by $0.03\%$ at
$h$, where Palace adds $4{,}946$ (initial) and $5{,}944$ (final) elements
by local bisection of interior boundary tetrahedra, and by $0.0000\%$ at
$h/2$, where no bisection occurs and both codes solve the same discrete
problem. Rerunning one cell with AMS in place of SuperLU, or feeding
geode-fem a format round-trip of the mesh, reproduces the same values.

\paragraph{Scope.}
The converged values are estimates for DeviceLayout's geometry under the
PEC-box, linear-junction model, from one mesh family per geometry refined
uniformly; the band does not include mesh-family spread. They are not a
statement about a fabricated device.

\section{A Static Route: Certified Island Capacitance}
\label{sec:static}

\paragraph{The certificate.}
For the reduced electrostatic problem (island at potential 1, ground and
resonator at 0, feedline floating with zero net charge), the potential
$\phi$ satisfies $0 \le \phi \le 1$ by the maximum principle, and
$C_\Sigma = \int \varepsilon\,\nabla\phi\cdot\nabla\phi$ with $\varepsilon
= \varepsilon_0\varepsilon_r$. Every conforming finite-element energy is an
upper bound on $C_\Sigma$, because $C_\Sigma$ is the minimum of the energy
over admissible potentials. For the lower bound, take any piecewise
polynomial field $D$ on the mesh cut along the zero-thickness sheets.
Integrating by parts and using $0\le\phi\le1$ on every term gives
\begin{equation}
  \int D\cdot\nabla\phi \;\ge\; N(D) \;=\; Q_{\mathrm{island}}(D)
  - \lvert Q_{\mathrm{feed}}(D)\rvert
  - \lVert \operatorname{div} D\rVert_{L^1}
  - \lVert [D\cdot n]\rVert_{L^1(\mathcal F)},
  \label{eq:cert-n}
\end{equation}
where $Q$ are the fluxes of $D$ into the two conductors and
$[D\cdot n]$ is the normal jump on the interior faces $\mathcal F$ off
the sheets. Cauchy--Schwarz, $\bigl(\int D\cdot\nabla\phi\bigr)^2 \le
\int D\cdot\varepsilon^{-1}D \int\varepsilon\nabla\phi\cdot\nabla\phi$,
then gives, whenever $N(D) > 0$,
\begin{equation}
  C_\Sigma \;\ge\; \frac{N(D)^2}{\int D\cdot\varepsilon^{-1}D}.
  \label{eq:cert}
\end{equation}
The bound holds for any $D$, so solver tolerance can only weaken it. We
take $D$ from the P2 potential, equilibrated patch by patch into
Raviart--Thomas fluxes of order one to three, with an independent
Crouzeix--Raviart flux as a check. The construction is validated against
an exact parallel plate, an analytic coaxial bracket and closed-form
checks of each term, with mutation tests that reject a grounded-feedline
or jump-free certificate. \emph{Certified} in this paper means rigorous
for the discrete polygonal geometry of the mesh, with lower bounds
floored and upper bounds ceiled. The certificate is evaluated in plain
f64, not interval arithmetic; the worst-case summation error is about
$10^{-10}$ relative, far below the rounding. Refinement does not move the
polygonal surfaces, so all meshes of one geometry bound the same
converged value, and a family bracket is the largest lower bound and the
smallest upper bound.

\paragraph{The unmodified device.}
Uniform refinement of the fixture with P1 and P2 elements gives the upper
bounds of Table~\ref{tab:convergence}. On adaptively refined meshes the
certificate closes the gap: the converged island capacitance of the
unmodified initial design $X^0$ lies in $[77.40, 77.61]$\,fF with the
scalar $\bar\varepsilon$ and in $[79.56, 79.76]$\,fF with the sapphire
tensor (Figure~\ref{fig:convergence}). The certified tensor-minus-scalar
difference is $1.95$--$2.35$\,fF, or $2.5$--$3.1\%$ of the scalar value,
and the same relative range holds for the two designs of
Appendix~\ref{sec:appendix-supporting}. The committed mesh overstates the
tensor $C_\Sigma$ by at least $56.92$\,fF, which is at least $71.3\%$ of
the converged value, and the scalar one by at least $77.4\%$.

\begin{table}[tbp]
  \centering
  \caption{The island capacitance $C_\Sigma$ (fF) of the unmodified device
  $X^0$ on uniform refinements of the committed mesh, and the certified
  brackets of its converged value from adapted meshes. Each uniform entry is
  a conforming energy and therefore an upper bound. Sources:
  \texttt{c\_sigma\_convergence.toml}, \texttt{c\_sigma\_bracket.toml} and
  \texttt{c\_sigma\_bracket\_tensor.toml} in
  \texttt{benchmarks/transmon\_diffopt/}.}
  \label{tab:convergence}
  \small
  \begin{tabular}{@{}lrrrrr@{}}
    \toprule
    Level & Tetrahedra & P1, scalar & P2, scalar & P1, tensor & P2, tensor \\
    \midrule
    $h$   & $133{,}314$     & $137.706752$ & $90.044141$ & $136.684731$ & $91.717846$ \\
    $h/2$ & $1{,}066{,}512$ & $101.359397$ & $83.015285$ & $102.375301$ & $85.008425$ \\
    $h/4$ & $8{,}532{,}096$ & $87.753422$  & n/a         & $89.493342$  & n/a \\
    \midrule
    \multicolumn{2}{@{}l}{Certified converged value} & \multicolumn{2}{c}{$[77.40, 77.61]$} & \multicolumn{2}{c}{$[79.56, 79.76]$} \\
    \bottomrule
  \end{tabular}
\end{table}

\begin{figure}[tbp]
  \centering
  \anvilfig{fig3-convergence.pdf}{fig3\_convergence.py}
  \caption{Mesh convergence of the island capacitance $C_\Sigma$ of the
  unmodified device under uniform refinement of the committed mesh
  (Table~\ref{tab:convergence}), with the scalar trace-averaged
  permittivity and with the sapphire tensor. Every marker is a conforming
  upper bound on the converged value of the committed polygonal geometry.
  The bars are the certified two-sided brackets from adapted meshes, drawn
  at the size of the mesh that sets them and magnified in the inset; the
  inset title gives the certified tensor-minus-scalar gap. Sources:
  \texttt{c\_sigma\_convergence.toml}, \texttt{c\_sigma\_bracket.toml}
  and \texttt{c\_sigma\_bracket\_tensor.toml}.}
  \label{fig:convergence}
\end{figure}

\paragraph{From capacitance to the qubit frequency.}
Reading the certified tensor bracket through the two-node circuit of
Section~\ref{sec:junction}, over every corner of the $C_\Sigma$, $L_s$ and
$k$ bands, with the resonator at its converged estimate, gives the
initial dressed qubit in $[4.449, 4.458]$\,GHz. That range lies inside the
eigenmode band of Section~\ref{sec:crosscheck}, and the eigenmode estimate
is $0.22\%$ below its midpoint. Because the extrapolation rule was chosen
after this range was known (Section~\ref{sec:crosscheck}), the $0.22\%$
is a consistency check, not an independent prediction. The range is an estimate, because $L_s$ and
$k$ are. One statement is sharper. A series inductance and a coupling to
a higher mode can only lower the qubit, so under the circuit model the
certified lower bound $C_\Sigma \ge 79.56$\,fF caps it at
$1/(2\pi\sqrt{L_J(C_\Sigma + C_J)}) = 4.477$\,GHz. The cap trims the high side of the
eigenmode band, which comes from extrapolation at the observed
pre-asymptotic order ($4.677$\,GHz). The cap does not depend on the
extrapolation rule, and the finest computed levels, Palace at order 2
on $h/2$, already lie $4.1\%$ and $4.7\%$ above the published initial
and final qubit values. Two
caveats apply: the bracket is certified on the fixture, which an
earlier generator produced from the same DeviceLayout geometry, not on
the meshes of the ladder; and the circuit is a model of the full-wave
problem.

\section{Gradient-Assisted Re-Optimization}
\label{sec:reopt}

The cross-check says that the published final design does not meet its
targets at converged resolution. We repeat DeviceLayout's optimization
twice: once against its own discretization, to compare cost with the
published run, and once against converged frequencies.

\paragraph{The capacitor\_length gradient.}
DeviceLayout keeps the junction fixed and builds the island upward from
it, so a change of $c$ translates every node above the island rigidly,
stretches the island band and re-centres the box. geode-fem turns this
motion into a node velocity $\partial X/\partial c$, checks on the mesh
that every metal outline node in a stretched or blended band lies on an
edge the motion maps onto itself ($3{,}146$ nodes on the initial mesh),
and evaluates Eq.~\eqref{eq:hf-eig} with the sapphire tensor for both
modes. The contraction takes $0.44$\,s. On the initial mesh it matches a
central difference of two re-solves of the moved mesh ($\pm0.5\,\mu$m) to
every printed digit: $-2.443227\times10^{-3}$\,GHz/$\mu$m for the qubit
and $-3.964\times10^{-6}$\,GHz/$\mu$m for the resonator.

\paragraph{The total\_length column.}
DeviceLayout re-lays the meander when $L$ changes, so no fixed-topology
node motion represents it. A transmission-line model holds
$f_r(L+\delta)$ constant, with $\delta = 216\,\mu$m from one finite
difference of two DeviceLayout meshes ($L = 5000$ and $5250\,\mu$m), and
takes $\partial f_q/\partial L = 0$. $\delta$ is weakly determined:
remesh noise of about $0.15\%$ in either $f_r$ moves it by about
$\pm175\,\mu$m (recomputed from the two committed frequencies), and a
second estimate gives about $363\,\mu$m. The model enters
only the Jacobian, and every iterate is evaluated on the true objective,
so a poorer Jacobian costs steps, not accuracy.

\paragraph{Newton on the high-fidelity objective.}
DeviceLayout re-meshes every design, so the first-order objective is not
smooth in the parameters: from $c = 620$ to $650\,\mu$m the remeshed finite
difference of the qubit frequency is $-5.23\times10^{-3}$\,GHz/$\mu$m,
against a fixed-topology derivative of $-2.28\times10^{-3}$. The iteration
therefore runs Newton directly on the high-fidelity objective $F$, in
DeviceLayout's variables $(x_1, x_2)$, in which the frequencies are close
to linear. The $c$ column of its Jacobian is geode-fem's first-order
derivative scaled by the measured fidelity ratio $S = F/f_h$ at the same
design. Each step costs one geode-fem first-order solve with its
gradient ($19$--$25$\,s) and one high-fidelity evaluation; the one-time
$\delta$ calibration is setup and is not counted.

\begin{table}[tbp]
  \centering
  \caption{Newton iterates against DeviceLayout's discretization (Palace on
  its second-order mesh with its configuration file). Step~0 is the
  initial design. err is DeviceLayout's measure against $[3.0, 4.0]$\,GHz;
  its criterion is $10^{-4}$. Source:
  \texttt{benchmarks/transmon\_devicelayout\_optimize/results.toml}.}
  \label{tab:newton-dl}
  \small
  \begin{tabular}{@{}rccccc@{}}
    \toprule
    Step & $c$ ($\mu$m) & $L$ ($\mu$m) & $f_q$ (GHz) & $f_r$ (GHz) & err \\
    \midrule
    0 & 620.000 & 5000.000 & 4.1683 & 5.6132 & $1.6\times10^{-1}$ \\
    1 & 1338.338 & 7137.638 & 2.8859 & 4.0465 & $7.9\times10^{-4}$ \\
    2 & 1225.346 & 7223.250 & 3.0120 & 4.0009 & $8.0\times10^{-6}$ \\
    3 & 1236.360 & 7224.900 & 3.0013 & 4.0011 & $1.4\times10^{-7}$ \\
    \bottomrule
  \end{tabular}
\end{table}

\paragraph{Against DeviceLayout's discretization.}
Table~\ref{tab:newton-dl} lists the iterates. The run meets DeviceLayout's
criterion at step~2, after 3 Palace runs counting the initial one as
DeviceLayout does, against PRIMA's 9; the three geode-fem solves add about
$70$\,s, against about $4$\,min per Palace run on 32 ranks
(Figure~\ref{fig:reopt}b). Step~3 reaches $\mathrm{err} =
1.4\times10^{-7}$ at $(1236.360, 7224.900)\,\mu$m, $+18.7$ and
$-19.1\,\mu$m from the published $(1217.678, 7243.955)\,\mu$m. About
$+10\,\mu$m of the $c$ offset is the Palace build: ours gives
$[3.0164, 3.9956]$\,GHz at the published final design, against the
published $[3.005, 3.982]$\,GHz. The rest is where PRIMA stopped; its
published pair sits $0.17\%$ and $0.45\%$ from the targets, inside its
$1\%$ rule. The comparison is with PRIMA's published run count, not with
its trajectory, which is not published.

\paragraph{Against the converged frequencies.}
Starting from the published final design, whose converged estimate is
$3.265$ and $4.088$\,GHz, each Newton step uses the converged estimate and the geode-fem
gradient at the current design, and each new design is evaluated by a
full ladder (Palace at order 1 on $h$, $h/2$, $h/4$ and order 2 on $h$,
$h/2$, geode-fem at order 1 on $h$, $h/2$) with the rule of
Section~\ref{sec:crosscheck}. Table~\ref{tab:newton-conv} gives the two
steps. The first overshoots the qubit by $-1.09\%$; the second lands at
$+0.16\%$ and $-0.04\%$ of the targets, with both targets inside the
bands. On both ladders geode-fem and Palace agree on the qubit to
$2\times10^{-4}\,\%$.

\begin{table}[tbp]
  \centering
  \caption{Newton steps against the converged frequencies. Each design is
  evaluated by a full resolution ladder; the converged values are
  estimates with extrapolation bands, not bounds. Source:
  \texttt{benchmarks/transmon\_devicelayout\_optimize/results.toml}
  and, for the published final design,
  \texttt{benchmarks/transmon\_devicelayout\_crosscheck/results.toml}.}
  \label{tab:newton-conv}
  \small
  \begin{tabular}{@{}lcccc@{}}
    \toprule
    Design & $c$ ($\mu$m) & $L$ ($\mu$m) & $f_q$ [band] (GHz) & $f_r$ [band] (GHz) \\
    \midrule
    Published final & 1217.678 & 7243.955 & 3.265 $[3.236, 3.471]$ & 4.088 $[4.077, 4.139]$ \\
    Step 1 & 1481.657 & 7407.945 & 2.967 $[2.938, 3.188]$ & 4.004 $[3.993, 4.055]$ \\
    Step 2 & 1446.078 & 7416.352 & 3.005 $[2.979, 3.221]$ & 3.998 $[3.987, 4.067]$ \\
    \bottomrule
  \end{tabular}
\end{table}

\paragraph{The converged design.}
At converged resolution the $3.0$/$4.0$\,GHz design needs a substantially
longer capacitor than the published one: the point estimate is
$c \approx 1446\,\mu$m, $+19\%$ ($+228\,\mu$m), with a resonator $2.4\%$
longer ($+172\,\mu$m). Mapped linearly through $dF/dc \approx
S\,\partial f_h/\partial c = -9.05\times10^{-4}$\,GHz/$\mu$m, the
slope at the step-1 design where step~2 started, and centred on the
step-2 estimate of $3.005$\,GHz, the step-2 qubit band corresponds to
$c$ between about $1418$ and $1685\,\mu$m, $+16\%$ to $+38\%$. The point estimate probably leans low:
the order-1 and order-2 series share a limit at order $0.87$ (qubit) and
$0.84$ (resonator), which puts step~2 at $3.039$ and $4.019$\,GHz,
$+1.3\%$ and $+0.5\%$ over the targets, so the converged design is
probably at or above the point-estimate $c$. Step~2 therefore lands on
the targets within the band; it is not a verified hit of $3.0$/$4.0$\,GHz
at convergence. Each ladder took $1.4$ and $1.6$\,h of Palace wall time,
dominated by the order-2 $h/2$ cell.

\begin{figure}[tbp]
  \centering
  \anvilfig{fig4-reoptimization.pdf}{fig4\_reoptimization.py}
  \caption{The gradient-assisted re-optimization. (a) The design plane near
  the optimum: the published final design, the Newton iterates against
  DeviceLayout's discretization (numbered; the first step starts from the
  initial design, off the plot), and the two steps against the converged
  frequencies, with the step-2 qubit band mapped linearly to
  capacitor\_length (an estimate). (b) DeviceLayout's err against the
  number of Palace runs at its discretization; PRIMA's published end point
  is plotted at its 9 runs, with err computed from the published
  frequencies, since its trajectory is not published. (c, d) Converged
  estimates and extrapolation bands of the published final design and the
  two converged-target steps, against the targets. Sources:
  \texttt{benchmarks/transmon\_devicelayout\_optimize/} and
  \texttt{benchmarks/transmon\_devicelayout\_crosscheck/results.toml}.}
  \label{fig:reopt}
\end{figure}

\section{The Sensitivity Layer}
\label{sec:adjoint}

For a solve $A(\varepsilon, X)\,x = b(X)$ and an observable $g(x)$, the
discrete adjoint $A^{\mathsf T}\lambda = \partial g/\partial x$ gives
\begin{equation}
  \frac{dg}{d\varepsilon_k}
    = -\lambda^{\mathsf T}\frac{\partial A}{\partial \varepsilon_k}x,
  \qquad
  \frac{dg}{dX}
    = \lambda^{\mathsf T}\frac{\partial b}{\partial X}
      - \lambda^{\mathsf T}\frac{\partial A}{\partial X}x ,
  \label{eq:adjoint}
\end{equation}
and in the complex driven case $2\,\mathrm{Re}[\cdot]$ of the same
contractions with the transpose solve, since the pencil is
complex-symmetric. The element Jacobians are exact, computed by
forward-mode dual numbers through the production kernels, and each
gradient costs one forward and one adjoint solve on one factorization.
Table~\ref{tab:adjoint-matrix} lists the four modules. Each has a test
against a central difference of the full pipeline (perturb, re-assemble,
re-solve, re-extract) and mutation tests: the conjugation error
$A^{\mathsf H}$ for $A^{\mathsf T}$ and a dropped term both fail. In the
driven shape case the source $b = \mathrm{i}\omega\mu_0\int N\cdot J$
depends on the geometry, and omitting $\partial b/\partial X$ fails the
check.

\begin{table}[tbp]
  \centering
  \caption{The discrete-adjoint sensitivity layer. Each cell: one forward
  and one adjoint solve on one factorization, exact element Jacobians, and
  the observed agreement with a central finite difference of the full
  pipeline. Modules in \texttt{crates/geode-core/src/}.}
  \label{tab:adjoint-matrix}
  \footnotesize
  \begin{tabular}{@{}llll@{}}
    \toprule
    Operator & Parameter & Module & FD rel-err \\
    \midrule
    scalar electrostatic (SPD) & material $\partial/\partial\varepsilon$ &
      \texttt{adjoint} & ${\sim}3\times10^{-8}$ \\
    scalar electrostatic (P1) & geometry $\partial/\partial X$ &
      \texttt{shape} & ${\sim}1\times10^{-9}$ \\
    driven $H(\mathrm{curl})$ & material $\partial/\partial\varepsilon$ &
      \texttt{driven::adjoint} & $2.3\times10^{-5}$ (worst region) \\
    driven $H(\mathrm{curl})$ & geometry $\partial/\partial X$ &
      \texttt{driven::shape} & $2.2\times10^{-9}$, $1.2\times10^{-8}$ (two maps) \\
    \bottomrule
  \end{tabular}
\end{table}

For the capacitance $C = \phi^{\mathsf T}K(X)\phi$ the stored energy is
stationary in the solved potential, so the implicit term vanishes and
$dC/dX = \phi^{\mathsf T}(\partial K/\partial X)\phi$, the classical
self-adjoint result \citep{park2019continuum}. The implementation keeps
the general adjoint form, and its test measures the implicit part below
$10^{-8}$ of the total. Composed with $E_C = e^2/(2(C_\Sigma + C_J))$,
this gives $\partial E_C/\partial X$. On a parallel-plate fixture Newton
reaches a target $E_C$ in one step, as expected for its affine
parameterization, and a fresh solve reproduces the target to
$1.4\times10^{-15}$. Appendix~\ref{sec:appendix-supporting} applies the
same identity, in a native-f64 P2 implementation outside the Burn tape,
to a three-parameter electrostatic optimization of the real device. The
eigenfrequency derivative of Section~\ref{sec:reopt} is the analogous
stationarity result for the eigenproblem.

\section{How geode-fem Was Designed and Built}
\label{sec:built}

geode-fem implements the whiteroom L1--L4 operator specification of
Crutcher Dunnavant \citep{dunnavant2026whiteroom}, and his design input
shaped the solver's operator layer. The code, the experiments and the
drafts of this paper were produced by frontier AI models working as
role-separated agents: Anthropic's Claude (Opus, Fable and Sonnet models,
through Claude Code) and OpenAI models (through Codex). The agents were
coordinated by Loom \citep{loom}, an issue, worktree and pull-request
workflow, and by a shared Squad room.

Every change passed an independent agent review that recomputed the
numbers from committed artifacts, reproduced runs and mutation-tested the
verifier pins, and several claims in this paper were corrected or dropped
because of those reviews. Three examples follow. The $4.14$\,GHz of the
SingleTransmon example was first read as a design target and set a
capacitance anchor; checked against the primary source, it is the
frequency of the unmodified initial geometry, and the designs aimed at it
are now supporting material (Appendix~\ref{sec:appendix-supporting}). A
capacitance comparison that mixed the scalar-$\bar\varepsilon$ and
sapphire-tensor models was replaced by certified brackets for each model
and a certified gap between them (Section~\ref{sec:static}). The agreement
of the static route with the eigenmode estimate was first presented as an
independent check; the extrapolation rule had been fixed after the static
route was known, so it is reported as a consistency check
(Section~\ref{sec:crosscheck}).

Counts from repository history, up to this revision: the paper and its
benchmark artifacts changed through 48 merged pull requests from
2026-07-14, and 82 of the 85 commits involved carry a Claude co-author
trailer (Opus 66, Fable 10, Sonnet 6). Commits made through Codex carry no
such trailer and are not counted. The human author directed the research,
made every scoping and claim decision, and takes responsibility for the
content.

\section{Reproducibility}
\label{sec:repro}

Every result in this paper is read from a committed, scripted
artifact.\footnote{The geode-fem source is public at
\url{https://github.com/rjwalters/geode-fem}. Archival DOI to be added at
deposit.}
% TODO(operator): Zenodo DOI
Table~\ref{tab:artifacts} lists them. The versions are geode-fem commit
\texttt{3174015} for the same-mesh credential and later commits named in
each artifact, Palace \texttt{fba6a5b}, DeviceLayout.jl v1.15.0 and
Burn 0.21. The DeviceLayout ladders ran on an AWS \texttt{r8g.8xlarge}
(Graviton4, 32 vCPU); the capacitance brackets ran on a local 28-core,
96\,GiB workstation. The Palace runs consume an MSH~2.2 conversion of each mesh,
because MFEM rejects DeviceLayout's multi-block MSH~4.1 file; the
conversion keeps every node and physical group and reproduces the
eigenvalues bit for bit. The figure scripts in \path{figures/src/} read
the artifacts directly.

\begin{table}[tbp]
  \centering
  \caption{Artifacts behind each result, relative to the repository root.
  The generated artifacts are pinned by regression tests.}
  \label{tab:artifacts}
  \footnotesize
  \begin{tabular}{@{}>{\raggedright\arraybackslash}p{0.30\linewidth}>{\raggedright\arraybackslash}p{0.64\linewidth}@{}}
    \toprule
    Result & Artifact \\
    \midrule
    Same-mesh agreement (Table~\ref{tab:agreement}) & \path{benchmarks/transmon_eigen/results.toml}; Palace cells under \path{reference/fixtures/transmon_palace/} \\
    Cross-check (Table~\ref{tab:crosscheck}, Figure~\ref{fig:crosscheck}) & \path{benchmarks/transmon_devicelayout_crosscheck/}; Palace cells under \path{reference/fixtures/transmon_palace/results_crosscheck_1092/}; meshes from \path{reference/julia/generate_transmon_param.jl} \\
    Capacitance convergence and brackets (Table~\ref{tab:convergence}, Figure~\ref{fig:convergence}) & \path{benchmarks/transmon_diffopt/c_sigma_convergence.toml}, \path{c_sigma_bracket.toml}, \path{c_sigma_bracket_tensor.toml}, \path{c_sigma_bracket_anchor_93p95.toml} \\
    Circuit terms $L_s$ and $k$ & \path{benchmarks/transmon_loop_inductance/results.toml}, \path{benchmarks/transmon_coupling/results.toml} \\
    Re-optimization (Tables~\ref{tab:newton-dl}--\ref{tab:newton-conv}, Figure~\ref{fig:reopt}) & \path{benchmarks/transmon_devicelayout_optimize/}; Palace cells under \path{reference/fixtures/transmon_palace/results_optimize_1092/} \\
    Electrostatic design loop (Appendix~\ref{sec:appendix-supporting}) & \path{benchmarks/transmon_diffopt/refined_optimization/} \\
    Performance (Appendix~\ref{sec:appendix-perf}) & \path{benchmarks/transmon_bench_cpu/}, \path{benchmarks/gpu_driven_scaling/} \\
    \bottomrule
  \end{tabular}
\end{table}

\section{Discussion}
\label{sec:discussion}

\paragraph{What is established, and how.}
Table~\ref{tab:tiers} sorts the results by the kind of evidence behind
them. The brackets are the only bounds. The converged frequencies and the
re-optimized design are extrapolations with stated rules and bands, and
the static route rests on circuit terms that are themselves
extrapolated. What stands without any extrapolation rule is the
certified cap and the finest computed levels, already $4.1\%$ and
$4.7\%$ above the published qubit values.

\begin{table}[tbp]
  \centering
  \caption{Evidence tiers. Certified: a guaranteed bound for the meshed
  polygonal geometry. Measured: a committed solve on a stated
  discretization. Estimated: an extrapolation with a stated rule, not a
  bound. Consistency: an agreement that is not independent evidence.}
  \label{tab:tiers}
  \footnotesize
  \begin{tabular}{@{}>{\raggedright\arraybackslash}p{0.15\linewidth}>{\raggedright\arraybackslash}p{0.79\linewidth}@{}}
    \toprule
    Tier & Results \\
    \midrule
    Certified & $C_\Sigma$ of $X^0$: $[77.40, 77.61]$\,fF (scalar), $[79.56, 79.76]$\,fF (tensor); tensor-minus-scalar gap $2.5$--$3.1\%$; the qubit cap of $4.477$\,GHz, conditional on the circuit model and certified on the fixture's polygonal geometry, not on the ladder meshes \\
    Measured & Same-mesh agreement (Table~\ref{tab:agreement}); Palace at DeviceLayout's discretization within $0.7\%$ of the published pairs; every ladder level, the finest already $4.1\%$ and $4.7\%$ above the published qubit values; same-mesh agreement on the ladders; the gradient check; 3 Palace runs to DeviceLayout's criterion \\
    Estimated & Converged frequencies and bands (Table~\ref{tab:crosscheck}); the static-route qubit range; $L_s$ and $k$; the converged design $c \approx 1446\,\mu$m and its $+16$ to $+38\%$ range \\
    Consistency & The two routes agree to $0.22\%$ on the initial qubit; the extrapolation rule was fixed after the static route was known, so this is not a blind prediction \\
    Not claimed & Agreement with measurement; a converged hit of $3.0$/$4.0$\,GHz; a global optimum; eigenmode-participation (EPR) derivatives; a speed advantage over Palace \\
    \bottomrule
  \end{tabular}
\end{table}

\paragraph{Threats and limitations.}
Agreement between two codes bounds implementation error, not shared
modeling choices, and the $17.49$\,GHz port-sheet mode is an example of
a shared artifact; the rank-1 shunt removes it in geode-fem, and Palace
still has it. The extrapolations assume a rate the sheet-edge singularity
sets and do not include mesh-family spread, and their high sides are
wide. The circuit route uses a bracket certified on a different mesh of
the same initial geometry, and $L_s$ and $k$ were extrapolated from two or
three discretizations. The re-optimization's transmission-line model
rests on one weakly determined $\delta$; it affects the step count, not
the iterates' evaluation. The PRIMA comparison counts runs only. The
electrostatic certificate is evaluated in f64 and its code is validated
by tests, not formally verified. The model is lossless and linear: the
junction is $L_J \parallel C_J$, the readout ports are open in the ladders,
and the metal is a set of zero-thickness sheets. Only one device is
studied. The eigenmode branch is not yet differentiated with respect to
participation ratios or Kerr couplings; an interior eigensolve that is
robust at the physical shift is the missing piece, and a projected
Jacobi--Davidson iteration \citep{liang2026adaptive} with the eigenpair
derivatives of \citet{nelson1976simplified} is the planned route.

\section{Conclusion}
\label{sec:conclusion}

The published frequencies of the DeviceLayout.jl SingleTransmon example
are reproduced by Palace at the example's own discretization and are
$7$--$9\%$ below the converged qubit frequencies that two independently
implemented solvers reach under refinement. A certified bracket on the
island capacitance, read through a circuit model, is consistent with that
shift and caps the initial qubit. A Newton iteration with a geode-fem eigenfrequency
gradient in its Jacobian reaches DeviceLayout's criterion in a third of
the published Palace runs, and at converged resolution it moves the
optimized design to a capacitor about $19\%$ longer, within a band of
$16$--$38\%$. For design examples that are meant to be copied, the
discretization is part of the result, and a solver that supplies
gradients and certified bounds makes it cheap to report converged values.

\paragraph{Acknowledgments.}
We thank Crutcher Dunnavant for the whiteroom specification and for his
design input to the solver's operator layer, and the Palace and
DeviceLayout.jl developers for publishing the SingleTransmon example and
its frequencies. The AI-orchestrated development of geode-fem and of this
paper is described in Section~\ref{sec:built}.
% TODO(operator): approve acknowledgment wording

\clearpage
\appendix
\section{Supporting Material: An Electrostatic Design Loop With a Certified Acceptance Test}
\label{sec:appendix-supporting}

This appendix reports an electrostatic optimization of the real device
whose interest is its acceptance test, not its target. The target, an
island capacitance of $93.954$\,fF, was obtained by reading the published
initial-geometry frequency of $4.14$\,GHz as a design target through the
linear junction circuit. Since $4.14$\,GHz is the frequency of the
unmodified initial design, the target has no design meaning; the runs
remain a measured example of how much a target reached on an
optimization mesh is worth after refinement.

\paragraph{Method.}
Ten cycles of adaptive refinement turn the committed mesh into a
$448{,}279$-tetrahedron optimization mesh. Three junction-pinned harmonic
morph fields change the island length, the pad width and the cutout gap.
The gradient is the stationary-energy shape derivative of
Section~\ref{sec:adjoint} for the floating-feedline $C_\Sigma$, evaluated
with P2 fields and exact quadrature on each affine tetrahedron, in native
f64 outside the Burn tape. The problem is underdetermined (three
parameters, one target), and the accepted steps are those that reduce the
freshly solved residual and keep the minimum element-volume ratio above
$0.25$. The design is scalar-$\bar\varepsilon$.

\paragraph{Result.}
The optimizer raises $C_\Sigma$ from $78.387918$ to $93.953668$\,fF in four
accepted steps; at the final design all three sensitivities agree with
central differences of fully re-solved geometries to $4.16\times10^{-7}$
(rounded up). Seven held-out refinement rounds, none of which feed the
optimizer, lower the conforming value to $92.998560$\,fF at
$3{,}388{,}197$ tetrahedra, and the certificate of
Section~\ref{sec:static} on the same meshes brackets the converged
capacitance of the design in $[92.75, 93.00]$\,fF with the scalar model
and $[95.34, 95.58]$\,fF with the sapphire tensor. The target reached on
the optimization mesh is lost by about $1\%$ under refinement, and the
certificate bounds the loss. A second run aimed past that bias
(target $95.112986$\,fF) gives $[93.89, 94.14]$\,fF scalar and
$[96.51, 96.75]$\,fF tensor. Palace with second-order elements on the
exported adapted mesh agrees with geode-fem's $C_\Sigma$ to
$2.5\times10^{-6}$, which tests the two assemblers on one discretization
and says nothing about the discretization error.

\paragraph{What it shows.}
A gradient loop converges to the target it is given on the mesh it is
given, and a finite-difference check or a fresh solve on that mesh says
nothing about the discretization error. The held-out conforming series
and the equilibrated-flux certificate together give a two-sided
acceptance test about $0.25$\,fF wide, independent of the optimizer. The
same certificate, with the tensor, supplies the bracket of
Section~\ref{sec:static}.

\section{Performance Context}
\label{sec:appendix-perf}

The timings below qualify the cost of the solvers; none measures the
re-optimization or ranks the two codes. The CUDA
backend (burn-cuda 0.21) is f32-only today: the underlying cubecl runtime
currently disables f64\footnote{An upstream constraint of the
burn/cubecl stack at version 0.21 \citep{burn}: the CUDA runtime reports
no f64 support, so the f64 code paths are not exercised on the GPU.} --- so every GPU result is
mixed-precision-qualified.

\paragraph{CPU eigenmodes.}
On one otherwise idle EC2 \texttt{m6i.4xlarge}, with every request
verified to return the six reference modes of Table~\ref{tab:agreement},
Table~\ref{tab:cpu} gives median wall seconds and measured core-seconds.
Two of the geode-fem requests are tuned using the reference spectrum, and
Palace solves about $17\%$ more unknowns from the same input mesh, so
these are workload comparisons, not kernel speedups. At eight wide Palace
without field output is fastest; at one wide the geode-fem routes take
less time, but geode-fem writes no fields and has no counterpart to
Palace's error estimator.

\begin{table}[tbp]
  \centering
  \caption{Mode-verified CPU timings on one m6i.4xlarge: median wall
  seconds / measured core-seconds. All rows contain the six reference
  modes of Table~\ref{tab:agreement}. Three repeats per cell except Palace
  one-rank Save=0 (one run). The two requests tuned on the reference
  spectrum are marked $\dagger$. Palace always requests six modes above
  $4.5$\,GHz. Source:
  \texttt{benchmarks/transmon\_bench\_cpu/results\_like\_for\_like\_m6i.toml}.}
  \label{tab:cpu}
  \small
  \begin{tabular}{@{}lrrr@{}}
    \toprule
    Request & Extra modes & 1 wide & 8 wide \\
    \midrule
    geode 30 at 4.5 GHz$^\dagger$ & 24 & 34.45 / 34.44 & 29.08 / 40.71 \\
    geode 6 at 20 GHz$^\dagger$ & 0 & 33.21 / 33.20 & 27.26 / 38.87 \\
    geode port-aware 8 at 4.5 GHz & 2 & 67.26 / 67.25 & 51.12 / 91.72 \\
    Palace, Save=6 & 0 & 186.84 / 186.82 & 33.35 / 266.66 \\
    Palace, Save=0 & 0 & 113.00 / 112.98 & 23.64 / 188.89 \\
    \bottomrule
  \end{tabular}
\end{table}

\paragraph{Driven solves and memory.}
On assembled first-order driven problems, geode-fem's AMS preconditioner
\citep{hiptmair-xu,tzaniovkolevpanayotsvassilevski2018parallel} reduces
one-core solve time relative to Jacobi by $3.4$--$4.2\times$ across
$25{,}695$--$897{,}650$ edges, but Palace's AMS is faster at scale
($53.3$ against $181.3$\,s at the largest size). On a lumped-port cube
family up to $25{,}695$ edges, geode-fem's GPU-f32 matrix-free solve is
slower than both assembled CPU configurations at every measured size and
reaches parity with CPU matrix-free only at the largest. A geode-fem
direct eigensolve at $1{,}157{,}564$ unknowns completed in $565.5$\,s at
$92.2$\,GB peak memory, on a different host from Palace's $423.12$\,s run;
the direct factorization's memory, not a controlled crossover, is what
that shows. The DeviceLayout ladders of Section~\ref{sec:crosscheck}
used Palace for the largest cells (up to $12.6$ million unknowns), and
geode-fem's direct solves at $h/2$ peaked at $94$ and $136$\,GB.

\clearpage

\bibliographystyle{plainnat}
\bibliography{refs}

\end{document}
